%!TEX TS-program = xelatex
%!TEX encoding = UTF-8 Unicode
% Awesome CV LaTeX Template for CV/Resume
%
% This template has been downloaded from:
% https://github.com/posquit0/Awesome-CV
%
% Author:
% Claud D. Park <posquit0.bj@gmail.com>
% http://www.posquit0.com
%
% Template license:
% CC BY-SA 4.0 (https://creativecommons.org/licenses/by-sa/4.0/)
%


%-------------------------------------------------------------------------------
% CONFIGURATIONS
%-------------------------------------------------------------------------------
% A4 paper size by default, use 'letterpaper' for US letter
\documentclass[11pt, a4paper]{awesome-cv-cn}

% Configure page margins with geometry
\geometry{left=1.4cm, top=.8cm, right=1.4cm, bottom=1.8cm, footskip=.5cm}

% Specify the location of the included fonts
\fontdir[fonts/]

% Color for highlights
% Awesome Colors: awesome-emerald, awesome-skyblue, awesome-red, awesome-pink, awesome-orange
%                 awesome-nephritis, awesome-concrete, awesome-darknight
\definecolor{pkured}{RGB}{140, 0, 0}
\colorlet{awesome}{pkured}
% Uncomment if you would like to specify your own color
% \definecolor{awesome}{HTML}{CA63A8}

% Colors for text
% Uncomment if you would like to specify your own color
% \definecolor{darktext}{HTML}{414141}
% \definecolor{text}{HTML}{333333}
% \definecolor{graytext}{HTML}{5D5D5D}
% \definecolor{lighttext}{HTML}{999999}

% Set false if you don't want to highlight section with awesome color
\setbool{acvSectionColorHighlight}{true}

% If you would like to change the social information separator from a pipe (|) to something else
\renewcommand{\acvHeaderSocialSep}{\quad\textbar\quad}


%-------------------------------------------------------------------------------
%	PERSONAL INFORMATION
%	Comment any of the lines below if they are not required
%-------------------------------------------------------------------------------
% Available options: circle|rectangle,edge/noedge,left/right
% \photo{./examples/profile.png}
\name{许科诺}
\position{高级工程师{\enskip\cdotp\enskip}华为}
%\address{Building N4, 101 Software Avenue, Nanjing, Jiangsu, 210012}

\mobile{(+86) 13269762233}
\email{kenuo.xu@foxmail.com}
\homepage{https://witty-me.github.io/}
\github{Witty-me}
\googlescholar{loX-G0gAAAAJ}{Google Scholar}
% \linkedin{posquit0}
% \gitlab{gitlab-id}
% \stackoverflow{SO-id}{SO-name}
% \twitter{@twit}
% \skype{skype-id}
% \reddit{reddit-id}
% \medium{madium-id}
% \googlescholar{googlescholar-id}{name-to-display}
%% \firstname and \lastname will be used
% \googlescholar{googlescholar-id}{}
% \extrainfo{extra informations}

% \quote{``Be the change that you want to see in the world."}


%-------------------------------------------------------------------------------
\begin{document}

% Print the header with above personal informations
% Give optional argument to change alignment(C: center, L: left, R: right)
\makecvheader

% Print the footer with 3 arguments(<left>, <center>, <right>)
% Leave any of these blank if they are not needed
\makecvfooter
  {\today}
  {许科诺~~~·~~~个人简历}
  {\thepage}


%-------------------------------------------------------------------------------
%	CV/RESUME CONTENT
%-------------------------------------------------------------------------------

\cvsection{工作经历}
\begin{cventries}
	
	%---------------------------------------------------------
	\cventry
	{高级工程师} % Degree (这里指职位)
	{华为} % Institution
	{中国 南京/上海} % Location
	{2025.08 - 至今} % Date(s)
	{
		\begin{cvitems} % Description(s) bullet points
			\item {就职于数据通信产品线—WLAN技术实验室，主要聚焦AI技术在室内无线网络与感知中的应用（实习期转正获评Excellent）。}
			\item {主导端到端 AI 算法架构的设计与开发，负责面向智慧安防与建筑场景的智能感知模型（基于 Wi-Fi CSI 等信号）的预研与落地。}
			\item {构建完整的数据闭环（采集、清洗与增强），设计并训练基于 CNN 与 Attention 机制的神经网络，完成感知任务模型。}
			\item {针对数据正负样本不平衡问题与跨环境泛化问题，设计基于自编码器的仅利用负样本的异常检测模型；并引入对抗性领域自适应（ADDA）算法，完成“实验室至现网”的跨域迁移。}
			\item {进行算法的端侧轻量化部署与性能评估，并在现网环境中完成实测，保障神经网络模型落地产品。}
		\end{cvitems}
	}
	
	\cventry
	{研究实习生} % Degree
	{微软亚洲研究院} % Institution
	{中国 上海} % Location
	{2022.12 - 2023.09} % Date(s)
	{
		\begin{cvitems} % Description(s) bullet points
			\item {就职于上海无线组；导师：邱锂力教授。}
			\item {进行大语言模型与计算机网络系统的交叉领域研究，探索大模型在底层系统优化中的应用潜力。}
			\item {设计并构建了基于 LLM 的网络算法参数智能调优框架，利用大模型的理解与推理能力，取代传统的人工调优规则。}
		\end{cvitems}
	}
	
	%---------------------------------------------------------
\end{cventries}


\cvsection{教育背景}
\begin{cventries}
	
	%---------------------------------------------------------
	\cventry
	{计算机科学 博士}
	{北京大学}
	{中国 北京}
	{2020.09 - 2025.06}
	{
		\begin{cvitems} % Description(s) bullet points
			\item {就读于软硬件协同体系结构（SOAR）研究组；导师：许辰人副教授。}
			\item {研究方向：移动计算、无线网络中的算法及其在物联网、机器人与人机交互领域的应用。}
			\item {设计了支持并发传输的可见光反向散射（Backscatter）通信系统，以实现低延迟通信目标。}
			\item {设计了以脉冲相机（Spike Camera）为接收端的可见光通信系统，实现了高数据传输速率与高动态范围。}
			\item {设计了以激光雷达为接收端的液晶基准标记（Fiducial Marker）系统，有效扩展了读取范围并提升了测距精度。}
			\item {更多相关研究工作请参阅“发表论文”部分。}
		\end{cvitems}
	}
	
	\cventry
	{计算机科学 学士} % Degree
	{北京大学} % Institution
	{中国 北京} % Location
	{2016.09 - 2020.06} % Date(s)
	{
		\begin{cvitems} % Description(s) bullet points
			\item {获北京大学“优秀毕业生”称号。}
		\end{cvitems}
	}
	
	%---------------------------------------------------------
\end{cventries}


\cvsection{论文发表}
\begin{cventries}
	
	%---------------------------------------------------------
	\cventry
	{\textbf{Kenuo Xu}, Bo Liang, Jingyu Li, Chenren Xu}
	{RetroLiDAR: A Liquid-crystal Fiducial Marker System for High-fidelity Perception of Embodied AI}
	{ACM SenSys}
	{2025}
	{
		\begin{cvitems} % Description(s) bullet points
			\item {一种利用激光雷达和液晶实现的远距离、高测距精度的基准标记系统，适用于机器人与虚拟现实领域。}
		\end{cvitems}
	}
	
	\cventry
	{Chenren Xu, \textbf{Kenuo Xu}, Lilei Feng, Bo Liang}
	{RetroV2X: A New V2X Paradigm with Visible Light Backscatter Networking}
	{Fundamental Research}
	{2023}
	{
		\begin{cvitems} % Description(s) bullet points
			\item {一套实用的基于可见光的车联网通信系统。}
		\end{cvitems}
	}
	
	\cventry
	{\textbf{Kenuo Xu}, Kexing Zhou, Chengxuan Zhu, Shanghang Zhang, Boxin Shi, Xiaoqiang Li, Tiejun Huang, Chenren Xu}
	{When Visible Light (Backscatter) Communication Meets Neuromorphic Cameras in V2X}
	{ACM HotMobile}
	{2023}
	{
		\begin{cvitems} % Description(s) bullet points
			\item {采用脉冲相机（Spike Camera）作为可见光通信接收端，适配移动场景下的数据传输。}
		\end{cvitems}
	}
	
	\cventry
	{\textbf{Kenuo Xu}, Chen Gong, Bo Liang, Yue Wu, Boya Di, Lingyang Song, Chenren Xu}
	{Low-Latency Visible Light Backscatter Networking with RetroMUMIMO}
	{ACM SenSys}
	{2022}
	{
		\begin{cvitems} % Description(s) bullet points
			\item {设计编解码算法支持 8 条并发的可见光反向散射链路，并将网络延迟降低了 92.0\%。}
		\end{cvitems}
	}
	
	\cventry
	{Jingyu Li, Qingwen Yang, \textbf{Kenuo Xu}, Yang Zhang, Chenren Xu}
	{EchoSight: Streamlining Bidirectional Virtual-physical Interaction with In-situ Optical Tethering}
	{ACM CHI}
	{2025}
	{
		\begin{cvitems} % Description(s) bullet points
			\item {一套面向 AR 眼镜的“即视即控”光学用户交互系统。}
		\end{cvitems}
	}
	
	\cventry
	{Yanxiang Wang, Yiran Shen, \textbf{Kenuo Xu}, Mahbub Hassan, Guangrong Zhao, Chenren Xu, Wen Hu}
	{High-Speed Passive Visible Light Communication with Event Cameras and Digital Micro-Mirror}
	{ACM SenSys}
	{2024}
	{
		\begin{cvitems} % Description(s) bullet points
			\item {利用事件相机（Event Camera）和数字微振镜（DMD），将无源可见光通信的吞吐量提升了 16 倍。}
		\end{cvitems}
	}
	
	\cventry
	{Zhiyuan He, Aashish Gottipati, Lili Qiu, Xufang Luo, \textbf{Kenuo Xu}, Yuqing Yang, Fanscis Y. Yan}
	{Designing Network Algorithms via Large Language Models}
	{ACM HotNets}
	{2024}
	{
		\begin{cvitems} % Description(s) bullet points
			\item {利用大语言模型设计并优化算法，提升计算机网络相关算法性能。}
		\end{cvitems}
	}
	
	\cventry
	{Chenren Xu, Purui Wang, Tuochao Chen, Yue Wu, \textbf{Kenuo Xu}, Xieyang Xu, Yang Shen, Junrui Yang, Guojun Chen, Guobin Shen}
	{VLID: Visible Light Backscatter System for Battery-free Internet-of-Things}
	{IEEE/ACM Transactions on Networking}
	{Accepted}
	{
		\begin{cvitems} % Description(s) bullet points
			\item {一套面向无电池物联网组网的端到端可见光反向散射解决方案。}
		\end{cvitems}
	}
	
	\cventry
	{Purui Wang, Lilei Feng, Guojun Chen, Chenren Xu, Yue Wu, \textbf{Kenuo Xu}, Guobin Shen, Kuntai Du, Gang Huang, Xuanzhe Liu}
	{Renovating road signs for infrastructure-to-vehicle networking: a visible light backscatter communication and networking approach}
	{ACM MobiCom}
	{2020}
	{
		\begin{cvitems} % Description(s) bullet points
			\item {利用可传达动态附加信息的反光路标，提升自动驾驶的可靠性。}
		\end{cvitems}
	}
	
	\cventry
	{Yue Wu, Purui Wang, \textbf{Kenuo Xu}, Lilei Feng, Chenren Xu}
	{Turboboosting Visible Light Backscatter Communication}
	{ACM SIGCOMM}
	{2020}
	{
		\begin{cvitems} % Description(s) bullet points
			\item {通过高级调制方案，将可见光反向散射（VLBC）的数据传输速率在原型系统中提升了 8 倍，在仿真环境中提升了 32 倍。}
		\end{cvitems}
	}
	
	%---------------------------------------------------------
\end{cventries}

\cvsection{荣誉奖项}

\begin{cvhonors}
	
	%---------------------------------------------------------
	\cvhonor
	{五四体育奖} % Award
	{北京大学} % Event
	{中国 北京} % Location
	{2024} % Date(s)
	
	%---------------------------------------------------------
	\cvhonor
	{二等奖} % Award
	{泛在智能感知技术创新应用大赛} % Event
	{中国 杭州} % Location
	{2024} % Date(s)
	
	%---------------------------------------------------------
	\cvhonor
	{三好学生} % Award
	{北京大学} % Event
	{中国 北京} % Location
	{2022} % Date(s)
	
	%---------------------------------------------------------
	\cvhonor
	{一等奖} % Award
	{未来网络科技创新大赛} % Event
	{中国 南京} % Location
	{2021} % Date(s)
	
	%---------------------------------------------------------
	\cvhonor
	{优秀毕业生} % Award
	{北京大学} % Event
	{中国 北京} % Location
	{2020} % Date(s)
	
	%---------------------------------------------------------
	\cvhonor
	{休斯顿校友会奖学金} % Award
	{北京大学} % Event
	{中国 北京} % Location
	{2019} % Date(s)
	
	%---------------------------------------------------------
	\cvhonor
	{三好学生} % Award
	{北京大学} % Event
	{中国 北京} % Location
	{2019} % Date(s)
	
	%---------------------------------------------------------
\end{cvhonors}

\cvsection{学术服务}
\begin{cventries}
	
	%---------------------------------------------------------
	\cventry
	{计算机网络（实验班）} % Title/Course
	{助教} % Role
	{北京大学} % Institution
	{2019、2020，2021-2024（参与）} % Date(s)
	{
		\begin{cvitems} % Description(s) bullet points
			\item {负责课程的日常组织与学生答疑工作。}
			\item {负责布置作业、主讲辅导课以及实验的评估与批改。}
			\item {参与随堂测试的题目设计与试卷批改。}
			\item {指导部分学生的课程期末研究项目。}
		\end{cvitems}
	}
	
	\cventry
	{Proceedings of the ACM on Interactive, Mobile, Wearable and Ubiquitous Technologies (IMWUT)} % Journal
	{期刊审稿人} % Role
	{} % Location
	{2021} % Date(s)
	{}	
	
	%---------------------------------------------------------
\end{cventries}

%\input{cv/education.tex}
%\input{cv/skills.tex}
%\input{cv/experience.tex}
%\input{cv/extracurricular.tex}
%\input{cv/honors.tex}
%\input{cv/presentation.tex}
%\input{cv/writing.tex}
%\input{cv/committees.tex}


%-------------------------------------------------------------------------------
\end{document}
